\documentclass[10pt]{article}

% Venue-neutral source. Replace the class and author block when a target venue
% is selected; keep the evidence hierarchy and claim boundaries intact.
\usepackage[margin=1in]{geometry}
\usepackage{amsmath}
\usepackage{amssymb}
\usepackage{array}
\usepackage{booktabs}
\usepackage{graphicx}
\usepackage{hyperref}
\usepackage[numbers,sort&compress]{natbib}
\usepackage{url}

\newcommand{\lcp}{LCP-style soft penalty}
\newcommand{\scppo}{SC-PPO 3.8 PID}
\newcommand{\heur}{Revised heuristic}
\newcommand{\jacc}{joint-acceleration}
\newif\ifsetpaperbibstyle
\setpaperbibstyletrue

\title{Policy Local Sensitivity Regularization for Smooth Humanoid Locomotion}
\author{Anonymous Authors}
\date{}

\begin{document}
\maketitle

\begin{abstract}
Smooth humanoid locomotion often uses action-rate or torque-rate reward
penalties, indirect regularizers that require task-specific scalar tuning. We
study policy local sensitivity: the amplification of observation perturbations
into action changes by the policy map. Under a shared Humanoid-Gym rough-terrain
protocol, we compare three interventions: SC-PPO as a hard policy-Jacobian
constraint, a locally adapted LCP-style soft Jacobian/Lipschitz penalty, and a
revised action-rate reward heuristic. In a five-seed Isaac Gym evaluation, the
soft penalty gives the strongest overall Isaac profile, with zero fall rate,
lowest velocity error, lowest action jitter, highest return, and lowest measured
sensitivity; the heuristic has the lowest joint acceleration. In matched
five-seed MuJoCo replay, the soft penalty retains the lowest action jitter; in
a bounded secondary trace analysis, it also has the lowest absolute
actuation-effort proxies, whereas the heuristic leads on mean joint
acceleration, energy per meter, and return. Controlled policy
perturbations and observation-noise stress show a consistent mechanism chain
from lower policy sensitivity to lower policy-output amplification and more
stable action streams, while downstream closed-loop metrics can still split.
Policy-local-sensitivity regularization is therefore best understood as a
policy-output smoothness mechanism rather than as a sufficient condition for
improving every closed-loop metric.
\end{abstract}

\section{Introduction}

Reinforcement-learning locomotion policies can satisfy task reward while
producing dynamically rough behavior. In humanoid locomotion this appears as
high-frequency action variation, large joint accelerations, and policies whose
smoothness depends on simulator details. A policy may perform well in its
training engine yet respond sharply to small observation or contact-solver
differences when replayed in another physics engine.

The standard engineering response is reward shaping. Action-rate, torque-rate,
joint-acceleration, and similar penalties are simple to implement and often
effective. However, the smoothness trade-off then becomes a scalar tuning
problem. Moreover, the policy is not directly constrained to be locally stable
as a map from observations to actions; it is penalized only after it has emitted
actions.

This paper evaluates whether policy local sensitivity, the local derivative of
policy output with respect to observations, provides an effective and
interpretable target for smooth humanoid control. If a policy strongly amplifies
small observation differences into action differences, then simulator-specific
contact and integration differences may appear as action jitter and dynamic
roughness. Regularizing this amplification can therefore improve policy-output
smoothness, although its effect on downstream closed-loop metrics depends on the
enforcement mechanism and the control stack.

We study this question through a mechanism comparison organized around three
levels: the policy map, the policy-output stream, and downstream closed-loop
dynamics. The evaluated methods
include a hard SC-PPO policy-Jacobian constraint, a soft LCP-style
Jacobian/Lipschitz penalty, and a revised action-rate reward heuristic. The
LCP-style penalty performs best among the sensitivity-based methods in the
present same-task evaluation, while the heuristic is stronger on downstream
MuJoCo joint acceleration and return. SC-PPO provides the hard-constraint
counterpart, allowing us to compare enforcement mechanisms directly under seed
variation, checkpoint selection, and cross-engine replay.

The contributions are:
\begin{enumerate}
  \item We formulate and empirically evaluate policy local sensitivity as a
  smooth-control mechanism for humanoid locomotion.
  \item We compare hard policy-sensitivity enforcement, soft
  LCP-style Jacobian/Lipschitz regularization, and a revised action-rate
  heuristic under the same five-seed Isaac Gym protocol.
  \item We perform matched five-seed MuJoCo replay and show a mechanism split:
  LCP-style regularization has the lowest action jitter and the lowest
  absolute actuation-effort proxies, while the heuristic has lower mean joint
  acceleration, lower energy per meter, and higher mean return.
  \item We report robustness analyses covering paired bootstrap and IQM
  uncertainty views, selected-vs-final checkpoint dependence, local LCP
  coefficient sensitivity, controlled policy perturbations, and
  observation-noise stress without retraining. Supplementary diagnostics add
  trace-based MuJoCo physical proxies, actuator-bandwidth stress, bounded
  second-setting terrain replays, and a bounded single-seed second-morphology
  check.
  \item The appendix records bounded OmniSafe PPO-Lag and local CPO-style
  feasibility checks as context for alternative constrained-RL
  implementations.
\end{enumerate}

\section{Related Work}

Constrained RL separates reward maximization from costs or limits. CPO
formulates policy optimization under explicit constraints~\citep{Achiam2017CPO},
while Safety Gym and PPO-Lagrangian practice provide common infrastructure for
cost-aware policy learning~\citep{Ray2019SafetyGym}. PID-Lagrangian methods
address multiplier overshoot and responsiveness in constrained deep
RL~\citep{Stooke2020PID}. OmniSafe packages many safe-RL algorithms behind a
common infrastructure~\citep{Ji2024OmniSafe}. Recent humanoid work also uses
constrained optimization to regulate physical costs such as
energy~\citep{Huang2026ECO}. This paper belongs to that cost/reward separation
lineage, but the cost placement is different: the SC-PPO cost is
actor-internal, constraining the policy's local sensitivity rather than an
environment-side outcome.

Smoothness in learned locomotion is commonly encouraged through reward terms on
action rate, torque rate, or acceleration. These methods are practical and form
strong reward-shaping baselines. A different family regularizes the policy map
itself. Spectral normalization constrains network Lipschitz
constants~\citep{Miyato2018SN}. LCP directly targets smooth humanoid
locomotion through a Lipschitz/Jacobian-style policy
regularizer~\citep{Chen2024LCP}, with adjacent work exploring spectral-normalization
variants for Lipschitz-constrained policies~\citep{Shin2025SNLCP}. We evaluate
a same-task LCP-style adaptation under Humanoid-Gym rough terrain; it is not an
official reproduction of the original LCP checkpoints, codebase, or evaluation
protocol.

Cross-engine replay provides an intermediate validation setting between
single-engine evaluation and hardware testing. Isaac Gym enables high-throughput
GPU simulation for robot learning~\citep{Makoviychuk2021IsaacGym}, while MuJoCo
uses a different physics engine and contact implementation~\citep{Todorov2012MuJoCo}.
Humanoid-Gym supplies the training scaffold and zero-shot transfer
context~\citep{Gu2024HumanoidGym}. We train and evaluate in Isaac Gym, then
replay selected checkpoints in aligned MuJoCo to test whether smoothness metrics
are robust to simulator changes. This sim-to-sim role is related to work on
robustness to dynamics mismatch, such as dynamics
randomization~\citep{Peng2018DynamicsRandomization}. We use MuJoCo replay as a
no-retraining sim-to-sim diagnostic for selected policy checkpoints.

\section{Methods}

\subsection{Shared PPO Task}

All primary methods use the same Humanoid-Gym rough-terrain locomotion task and
PPO backbone~\citep{Schulman2017PPO}. The robot is an H1-class humanoid with
12 DoF. The main evaluation uses seeds 11, 17, 23, 29, and 31, a
512-environment training budget, and a checkpoint sweep over the trained policy.
Evaluation reports fall rate, velocity-tracking error, joint-acceleration
$L_2$, action-jitter $L_2$, episode return, and policy-sensitivity measurements
where available.

The selected-checkpoint protocol is fixed before comparison: checkpoints are
filtered by a task floor and then selected by the lowest joint-acceleration
$L_2$, with action jitter and tracking error as tie-breakers. Final-checkpoint
behavior is reported separately in the checkpoint robustness analysis.

\subsection{SC-PPO Hard Policy-Sensitivity Constraint}

SC-PPO introduces a constraint cost based on policy local sensitivity. For a
policy mean $\mu_\theta(o)$, the local cost is the Frobenius norm of the
observation-to-action Jacobian:
\begin{equation}
  C(o) = \left\| \frac{\partial \mu_\theta(o)}{\partial o} \right\|_F .
\end{equation}
The batch cost for multiplier updates is a tail-sensitive statistic:
\begin{equation}
  J_C(\theta) = Q_{0.90}\left(\{C(o_i)\}\right).
\end{equation}
SC-PPO optimizes a Lagrangian objective
\begin{equation}
  \mathcal{L}(\theta,\lambda) =
  J_R(\theta) - \lambda \left(J_C(\theta) - d\right),
\end{equation}
with threshold $d=3.8$, $\lambda_{\mathrm{init}}=0.5$, quantile aggregation,
and \texttt{lower\_bound\_clamp} PID integral handling. Heuristic smoothness
rewards are disabled in the SC-PPO condition, so smoothness pressure comes from
the Jacobian constraint rather than from reward shaping or any auxiliary
smoothness penalty.

\subsection{LCP-Style Soft Jacobian/Lipschitz Penalty}

The LCP-style method uses a fixed soft penalty on the policy-gradient quantity
$\|\nabla_o \log \pi_\theta(a \mid o)\|^2$, following the mechanism family of
Lipschitz-constrained policies~\citep{Chen2024LCP}. The local implementation
computes this penalty on PPO minibatch state-action samples with
\texttt{create\_graph=True}, uses a random subsample of actor observations for
computational tractability, and applies \texttt{lcp\_weight = 0.002}. It does
not use a Lagrange multiplier and disables heuristic smoothness rewards. It
shares the same task, checkpoint sweep, metric schema, and MuJoCo replay bridge
as SC-PPO and the heuristic.

This score-gradient penalty is related to, but not identical to, the SC-PPO
mean-Jacobian cost. For the diagonal Gaussian policy with state-independent
diagonal exploration standard deviation vector $\boldsymbol{\sigma}$ used here,
\begin{equation}
  \nabla_o \log \pi_\theta(a \mid o)
  =
  (a-\mu_\theta(o))^\top
  \operatorname{diag}(\boldsymbol{\sigma}^{-2})
  \frac{\partial \mu_\theta(o)}{\partial o}.
\end{equation}
Thus the LCP-style penalty is modulated by both the sampled action residual and
the current exploration scale. In the scalar-variance, fixed-residual
simplification, the squared term carries a $\sigma^{-4}$ factor; on-policy
residual variance partially offsets this factor, but the effective pressure
remains tied to exploration variance. Consequently, the effective
regularization pressure is weaker during early high-exploration phases and
stronger after exploration shrinks, giving the soft penalty a built-in
training-time schedule absent from SC-PPO's fixed mean-Jacobian threshold.

This condition is the closest literature-aligned policy-sensitivity baseline in
our experimental stack, but it is not an official LCP reproduction. The public
LCP/MimicKit stack uses a different task, robot, checkpoint setup, and
evaluation path; we therefore label this method LCP-style rather than claiming
official parity.
Throughout the paper, ``LCP-style'' denotes this locally adapted soft
Jacobian/Lipschitz penalty, not an official LCP reproduction; when context is
clear, we refer to the same row as the soft penalty.
Appendix Table~\ref{tab:lcp-parity} makes this boundary explicit. The current
row matches the soft penalty family, fixed-coefficient actor-update path, and
same-task comparison protocol used in this paper, but it does not match the
official task, robot, checkpoint/evaluation path, or hardware stack. We
therefore use it as a same-task mechanism baseline rather than as an
official-parity comparison.

\subsection{Revised Heuristic Reward-Shaping Anchor}

The heuristic method is PPO with a tuned action-rate penalty under the formal
protocol revision. It represents the practical reward-shaping approach that the
policy-sensitivity methods must compare against. Because it is strong in MuJoCo,
especially on joint acceleration and return, it is treated as a primary
baseline. The selected anchor uses
\texttt{rewards.scales.action\_smoothness = -0.0050}; the preceding
action-rate sweep tested $-0.0005$, $-0.0020$, and $-0.0050$ while holding
\texttt{dof\_acc = -1e-7}, \texttt{base\_acc = 0.2}, and
\texttt{dof\_vel = -5e-4} fixed. The same task-floor-then-smoothness
checkpoint-selection rule is applied.

\section{Experimental Protocol}

The primary Isaac comparison uses five seeds, checkpoint sweeps, and the same
metric schema for LCP-style, SC-PPO, and heuristic methods. MuJoCo replay uses the
selected Isaac checkpoint for each seed under the aligned \texttt{isaac\_mainline}
protocol with 20 episodes of 20 seconds and \texttt{joint\_reset\_noise = 0.1}.
For the selected MuJoCo replay we also compute secondary torque- and
power-based proxies from exported joint-control and joint-velocity traces.
Each selected checkpoint is rerun for five trace-capture episodes.
Absolute power is the mean absolute joint power, absolute episode energy
integrates absolute joint power over time, and energy per meter divides that
absolute energy by forward distance traveled; all are lower-is-better
descriptive proxies. These trace metrics are physically grounded supplementary
evidence; they do not replace the primary task, \jacc, or action-jitter
metrics.

Metrics are fall rate, velocity-tracking error, joint-acceleration $L_2$,
action-jitter $L_2$, episode return, and policy sensitivity. Fall, velocity
error, joint acceleration, action jitter, and sensitivity are lower-is-better;
episode return is higher-is-better. Policy sensitivity and violation rate are
auxiliary measurements of policy-local behavior rather than task reward.
Episode return is not normalized across engines: MuJoCo reports an
evaluator-local tracking proxy accumulated as forward velocity minus tracking
error, not the Humanoid-Gym training reward. Isaac and MuJoCo return values
should therefore be compared only within the same engine.

The primary comparison consists of the LCP-style soft penalty, SC-PPO 3.8
PID-Lagrangian, and the revised heuristic. Robustness analyses include paired
bootstrap uncertainty, selected-vs-final checkpoint dependence, and LCP
coefficient sensitivity, controlled policy perturbations, and no-retraining
observation noise. Supplementary diagnostics include no-retraining MuJoCo
actuator-bandwidth stress, \texttt{hfield\_moderate} replay, a bounded
same-morphology command shift, a bounded single-seed second-morphology check,
and bounded OmniSafe PPO-Lag and local CPO-style feasibility notes. These
additional slices are used to sharpen claim boundaries rather than to expand
the main comparison table.

\section{Results}

\subsection{Five-Seed Isaac Evaluation}

Table~\ref{tab:isaac} shows the primary five-seed Isaac Gym
selected-checkpoint comparison. The soft penalty achieves the best overall
Isaac results: zero fall rate, lowest velocity error, lowest action jitter,
highest return, and lowest sensitivity. The heuristic has the lowest joint
acceleration, while the hard-constraint SC-PPO row is weaker under this
enforcement setting.

\begin{table}[t]
\centering
\caption{Five-seed Isaac Gym selected-checkpoint comparison. Values are
mean $\pm$ standard deviation over seeds 11/17/23/29/31. Abbreviations:
Vel. err is velocity-tracking error, Jnt acc is joint-acceleration $L_2$,
Jitter is action-jitter $L_2$, and Sens. is policy sensitivity.}
\label{tab:isaac}
\small
\resizebox{\linewidth}{!}{%
\begin{tabular}{lccccccc}
\toprule
Method & Checkpoints & Fall & Vel. err & Jnt acc & Jitter & Return & Sens. \\
\midrule
\lcp & 300/400/400/400/400 & $0.000\pm0.000$ & $0.490\pm0.032$ & $117.331\pm21.171$ & $0.212\pm0.033$ & $118.420\pm6.595$ & $1.890\pm0.064$ \\
\scppo & 300/300/400/400/400 & $0.170\pm0.194$ & $0.606\pm0.087$ & $142.955\pm52.454$ & $0.277\pm0.100$ & $99.349\pm13.959$ & $3.630\pm0.067$ \\
\heur & 350/300/350/400/400 & $0.150\pm0.063$ & $0.705\pm0.103$ & $115.317\pm8.770$ & $0.260\pm0.018$ & $105.326\pm10.335$ & $7.331\pm0.237$ \\
\bottomrule
\end{tabular}}
\end{table}

A descriptive five-seed audit that combines paired bootstrap intervals, IQM,
and bootstrap rank stability preserves the same read as
Table~\ref{tab:isaac}: LCP is stronger than SC-PPO on Isaac fall rate,
velocity error, return, and measured sensitivity. The joint-acceleration
difference relative to SC-PPO is favorable in the mean but not resolved under
five seeds because the confidence interval overlaps zero. Relative to the
heuristic, LCP improves Isaac fall rate, velocity error, jitter, return, and
sensitivity, while the heuristic retains the lowest joint acceleration. All
three selected rows remain task-valid on all five seeds, so the primary
reliability distinction is checkpoint dependence rather than widespread
collapse.

\subsection{Matched MuJoCo Replay}

Table~\ref{tab:mujoco} reports matched five-seed MuJoCo selected replay. The
soft penalty has the lowest action jitter and improves substantially over
SC-PPO, while the revised heuristic achieves lower mean joint acceleration and
higher return.

\begin{table}[t]
\centering
\caption{Matched five-seed MuJoCo selected-checkpoint replay. Values are
mean $\pm$ standard deviation over the same selected checkpoints as
Table~\ref{tab:isaac}.}
\label{tab:mujoco}
\small
\resizebox{\linewidth}{!}{%
\begin{tabular}{lccccc}
\toprule
Method & Fall & Vel. err & Jnt acc & Jitter & Return \\
\midrule
\lcp & $0.000\pm0.000$ & $0.406\pm0.060$ & $117.425\pm11.234$ & $0.195\pm0.029$ & $-599.108\pm236.114$ \\
\scppo & $0.010\pm0.020$ & $0.471\pm0.079$ & $159.718\pm71.345$ & $0.322\pm0.167$ & $-627.238\pm222.350$ \\
\heur & $0.000\pm0.000$ & $0.406\pm0.043$ & $111.615\pm14.597$ & $0.226\pm0.037$ & $-456.370\pm195.030$ \\
\bottomrule
\end{tabular}}
\end{table}

Secondary trace-based physical proxies sharpen this split rather than reversing
it. On the same selected checkpoints, the soft penalty has the lowest torque
RMS ($24.102$), absolute power ($85.659$), and absolute episode energy
($877.146$), whereas the heuristic has the lowest energy-per-meter proxy
($785.610$) and the lowest aggregate MuJoCo joint acceleration; SC-PPO is worst
on absolute power ($184.312$) and aggregate joint acceleration
(Appendix Table~\ref{tab:mujoco-physical}). This result separates
policy-output smoothness and absolute actuation effort from downstream dynamic
and efficiency-style metrics. The sensitivity regularizer acts directly on the
policy-output stream, which is consistent with the soft penalty achieving the
lowest mean action jitter and the best action jitter on three of five seeds.
Joint acceleration, energy per meter, and return are downstream closed-loop
outcomes that also depend on velocity tracking, contact timing, PD target
dynamics, and simulator-specific response. A leave-one-seed analysis preserves
this metric split: the soft penalty remains the lowest mean action-jitter
method in every split, while the heuristic remains best on mean joint
acceleration and mean return in every split. Because this trace slice uses only
five captured episodes per selected checkpoint and is not bootstrap-backed, we
use it as bounded secondary evidence rather than as a new primary comparison.

\subsection{Mechanism-Chain Analysis}

Figure~\ref{fig:mechanism-chain} is the central interpretive result of the
paper. It organizes the comparison into three linked but distinct levels: the
policy map, the policy-output stream, and downstream closed-loop dynamics. The
main claim is that the soft penalty is consistently strongest on the first two
levels, while the third can still split in favor of the heuristic on selected
metrics.

A controlled policy perturbation analysis provides a direct measurement of
local policy-output amplification. On a shared observation bank from
selected-checkpoint rollouts, deterministic action responses to small
observation perturbations follow the expected ordering: the soft penalty has
the lowest amplification, SC-PPO is intermediate, and the revised heuristic is
largest.

\begin{table}[t]
\centering
\caption{Policy perturbation analysis on a shared observation bank. Amplification
is $\|\Delta a\|/\|\Delta o\|$ at perturbation radius $0.01$ using deterministic
action means.}
\label{tab:perturb}
\small
\resizebox{\linewidth}{!}{%
\begin{tabular}{lccccccc}
\toprule
Method & Seeds & Amp. mean & Amp. p90 & Bank sens. & Selected sens. & Jitter & Jnt acc \\
\midrule
\lcp & 5 & 0.082 & 0.133 & 2.308 & 1.890 & 0.212 & 117.331 \\
\scppo & 5 & 0.153 & 0.241 & 4.307 & 3.630 & 0.277 & 142.955 \\
\heur & 5 & 0.289 & 0.468 & 8.341 & 7.331 & 0.260 & 115.317 \\
\bottomrule
\end{tabular}}
\end{table}

Across the 15 method-seed conditions, amplification mean correlates with
selected policy sensitivity at $0.999$. This near-identity validates the
sensitivity construct: two independent measurement methods converge to the same
ordering across all method-seed conditions. The weaker correlations with
selected action jitter ($0.228$) and selected joint acceleration ($-0.093$) show
why downstream closed-loop metrics should not be treated as direct sensitivity
measurements.

The policy-map and policy-output measurements align: the soft penalty has the
lowest measured sensitivity, the smallest response to controlled observation
perturbations, the lowest matched MuJoCo action jitter, and the smallest
action-jitter degradation under observation noise. The downstream MuJoCo
joint-acceleration and return reads can split from this ordering, and the
secondary physical-proxy table shows that energy per meter can split in the
same downstream stage. The figure is therefore the main interpretive summary
of the mechanism comparison, not a causal graph or a hardware-transfer claim.

\begin{figure}[t]
\centering
\includegraphics[width=\linewidth]{../../artifacts/analysis/paper_figures/figure_mechanism_chain.png}
\caption{Mechanism-chain evidence for selected checkpoints. Columns are grouped
into policy-map, policy-output, and downstream closed-loop evidence; rank 1 is
best within each column. The policy-map and policy-output levels follow the
soft-penalty $<$ SC-PPO $<$ heuristic ordering, while the downstream MuJoCo
joint-acceleration and return columns favor the heuristic in the matched replay
of Table~\ref{tab:mujoco}. The observation-noise column reports the worse of
Isaac and MuJoCo jitter factors at $\sigma=0.05$.}
\label{fig:mechanism-chain}
\end{figure}

\paragraph{Observation-noise stress.}
We next test whether the local mechanism survives a closed-loop perturbation
without retraining. Gaussian noise is added only to the policy observation
input at inference time, using the same selected checkpoints and the same five
seeds. Table~\ref{tab:obs-noise} reports the largest tested noise scale,
$\sigma=0.05$; the full sweep, omitted from the table for space, also includes
$0.005$, $0.01$, and $0.02$.

\begin{table}[t]
\centering
\caption{No-retraining observation-noise robustness at $\sigma=0.05$. Jitter
factor is action-jitter $L_2$ relative to the corresponding no-noise replay in
the same engine.}
\label{tab:obs-noise}
\small
\resizebox{\linewidth}{!}{%
\begin{tabular}{lcccc}
\toprule
Method & Isaac fall & Isaac jitter factor & MuJoCo fall & MuJoCo jitter factor \\
\midrule
\lcp & 0.000 & 1.095 & 0.000 & 1.238 \\
\scppo & 0.250 & 1.219 & 0.000 & 1.401 \\
\heur & 0.260 & 1.800 & 0.000 & 2.955 \\
\bottomrule
\end{tabular}}
\end{table}

The degradation ordering follows the policy-sensitivity and perturbation-audit
ordering: the soft penalty has the smallest action-jitter increase in both
engines, SC-PPO is intermediate, and the heuristic is most affected. This
strengthens the policy-output robustness interpretation. It remains an
input-perturbation stress test rather than hardware validation, and it does not
override the downstream metric split observed in matched MuJoCo replay.

\subsection{Robustness Analyses}

The statistical robustness analysis is descriptive rather than a large-sample
significance test. With five seeds, we report interval estimates, IQM, and
rank-stability views to characterize effect directions and uncertainty rather
than binary accept/reject conclusions. Across paired bootstrap intervals, IQM,
and bootstrap rank stability, the most conservative robust comparison is that
LCP has stable advantages over the SC-PPO hard-constraint condition on Isaac
fall rate, velocity error, return, sensitivity, and MuJoCo action jitter.
Against the heuristic, the same robust aggregate preserves a trade-off rather
than a dominance claim: LCP is better on Isaac task behavior, sensitivity, and
action jitter, whereas the heuristic remains competitive or better on joint
acceleration, especially in MuJoCo.

The selected-vs-final checkpoint analysis shows that LCP is close to final-only
behavior. Under the selected-checkpoint protocol, all three methods remain
task-valid on all five seeds (\texttt{fall\_rate} $< 1.0$), so the main
reliability split is checkpoint dependence rather than widespread collapse.
Only seed 11 selects checkpoint 300 instead of final checkpoint 400 for LCP,
giving selected-final agreement on $4/5$ seeds; SC-PPO and the heuristic agree
with final on $3/5$ and $2/5$ seeds, respectively. SC-PPO is more
dynamic-selection-sensitive: final checkpoints improve velocity and return but
worsen joint acceleration and jitter. The heuristic is task-selection-sensitive:
final checkpoints improve velocity and return slightly but increase fall rate
and dynamic roughness. Table~\ref{tab:checkpoint-robustness} gives the compact
delta read.

\begin{table}[t]
\centering
\caption{Selected-vs-final checkpoint robustness. Deltas are final minus
selected over seeds 11/17/23/29/31; for lower-is-better metrics, positive
values mean the final checkpoint is worse than the selected checkpoint.}
\label{tab:checkpoint-robustness}
\small
\resizebox{\linewidth}{!}{%
\begin{tabular}{lcccccc}
\toprule
Method & Isaac Fall & Isaac Jnt acc & Isaac Jitter & MuJoCo Fall & MuJoCo Jnt acc & MuJoCo Jitter \\
\midrule
\lcp & $0.000$ & $1.004$ & $0.007$ & $0.000$ & $-0.170$ & $0.008$ \\
\scppo & $0.000$ & $8.932$ & $0.014$ & $-0.010$ & $8.689$ & $0.016$ \\
\heur & $0.050$ & $2.463$ & $0.009$ & $0.000$ & $0.249$ & $-0.001$ \\
\bottomrule
\end{tabular}}
\end{table}

The LCP coefficient analysis tests weights 0.001, 0.002, and 0.004 on the same
five-seed set. The smallest weight appears under-regularized, with nonzero fall
rate, higher measured sensitivity, and rougher selected policies. The largest
weight lowers measured sensitivity further but worsens velocity error, joint
acceleration, and action jitter relative to 0.002, while return is nearly tied.
This does not prove global hyperparameter optimality, but it shows that the main
LCP condition is not an isolated single-point accident within the tested
neighborhood.

Supplementary robustness slices are reported in the appendix rather than
promoted into the main comparison story. These include the no-retraining
\texttt{hfield\_moderate} replay, which remains useful as a bounded
terrain-generality check but operates at materially higher fall rates than the
matched \texttt{isaac\_mainline} MuJoCo replay, and the actuator-bandwidth
stress test, whose $\tau=0.05$ regime degrades all methods almost
identically while $\tau=0.10$ mostly exposes a task-validity boundary. Two
bounded terrain-aware main-morphology retrain probes also remained
task-invalid: a stair-bearing mixed rough/stairs slice and a no-stairs
\texttt{trimesh + curriculum} slice both yielded \texttt{fall\_rate = 1.0}
across their evaluated checkpoints. We therefore treat terrain-aware retraining
as a negative-result boundary rather than as supplementary robustness evidence
in the current paper.

The SC-PPO threshold is a local design choice rather than a globally optimized
constant. A threshold sweep over nearby shared values found seed-dependent
behavior: lower settings over-constrained some seeds, while looser settings
improved one seed but degraded others. No tested shared threshold in that
local grid passed all task-validity and smoothness gates. We therefore use
threshold 3.8 as the documented hard-constraint setting. A matched plain-dual
comparison at the same threshold also failed the task-validity criteria,
supporting PID-Lagrangian as the stronger within-family SC-PPO condition without
establishing the independent necessity of every PID term.

\section{Discussion}

\subsection{What the Mechanism Comparison Shows}

The main conclusion is a three-level mechanism-chain decomposition. The soft
penalty most consistently improves policy-map sensitivity, controlled
perturbation amplification, action jitter, and observation-noise degradation.
The heuristic remains stronger on some downstream closed-loop metrics,
especially MuJoCo joint acceleration, energy per meter, and return. The
five-seed uncertainty audit therefore supports a mechanism-level claim rather
than a blanket statement that one method uniformly leads on every metric.

The evidence for this decomposition is internally consistent. Perturbation
amplification and selected policy sensitivity correlate at $0.999$, showing
that two independent measurements recover the same policy-output amplification
axis. In a bounded secondary trace analysis of the matched MuJoCo selected
checkpoints, the soft penalty is also lowest on torque RMS, absolute power, and
absolute episode energy, whereas the heuristic is best on energy per meter and
aggregate joint acceleration. Policy-local-sensitivity regularization therefore
appears most effective at reducing absolute control effort and output
roughness, while efficiency-style progress-normalized metrics and full
closed-loop dynamics can still favor a reward-shaped policy.

One plausible optimization explanation is that soft regularization is easier to
fit in PPO because the penalty is present throughout actor updates with a fixed
coefficient, and because the Gaussian score-gradient penalty is modulated by
the action residual and exploration scale rather than imposing a fixed hard
threshold on the mean Jacobian. SC-PPO exposes the behavior of hard
enforcement under the same policy-sensitivity objective: the constraint is
interpretable, but PID-Lagrangian multiplier dynamics and checkpoint behavior
are seed-sensitive. In this comparison, SC-PPO is therefore best read as a
hard-enforcement diagnostic rather than as the main performance reference.

For practitioners, the results suggest that adding a fixed soft Jacobian
penalty can be a simple alternative to reward-only smoothness tuning in a PPO
locomotion stack when policy-output smoothness and robustness are the primary
goals. If downstream joint acceleration, energy efficiency, or return are the
main priority, the revised heuristic remains a strong practical baseline.

A bounded single-seed Unitree H1 diagnostic under the same shared metric schema
provides a complementary warning. On this second morphology, the revised
heuristic preserved the repaired H1 task floor and improved downstream
smoothness relative to the H1 vanilla baseline, whereas the current LCP-style
and SC-PPO Jacobian-based mechanisms did not (Appendix
Table~\ref{tab:h1-diagnostic}). This supplementary result reinforces the main
paper's scope boundary: lowering the policy-sensitivity metric alone is not
sufficient for good downstream locomotion, and mechanism transfer can vary
substantially across morphologies.

A task-valid no-retraining command-shift slice on the same replay path provides
the complementary bounded-positive supplementary read. On a MuJoCo plane replay
with \texttt{command\_vx = 0.6}, the same metric split persists: LCP remains
best on action jitter, while the revised heuristic remains best on velocity
error, joint acceleration, and return (Appendix
Table~\ref{tab:command-shift}). This is useful supplementary robustness
evidence, but it is a command-distribution shift rather than a terrain
benchmark.

\subsection{Why the Heuristic Still Matters}

The revised heuristic is a strong reward-shaping baseline because it retains
advantages on some
downstream metrics without reducing policy sensitivity. In matched MuJoCo
replay it outperforms LCP on mean joint acceleration, energy per meter, and
mean return even though LCP achieves lower action jitter and lower absolute
effort proxies. This trade-off matters because it shows that reward shaping can
exploit full closed-loop dynamics in ways that are not captured by policy-map
smoothness alone.

\subsection{Sim-to-Sim Scope}

MuJoCo replay changes physics-engine details without changing the trained
policy checkpoint. The sim-to-sim result therefore tests whether smoothness
metrics are robust to simulator changes, but it does not establish sim-to-real
robustness. This matters because a smoothness mechanism with purely
simulator-specific benefits has limited practical value; the preservation of the
action-jitter ordering across engines, together with the perturbation and
observation-noise audits, shows that the soft penalty's policy-output
smoothness is not merely an Isaac Gym artifact. Domain randomization, system
identification, actuator modeling, and hardware tests remain outside the
present scope.

\section{Limitations}

The LCP-style method is a local same-task adaptation, not an official
reproduction of the original LCP code, checkpoints, or evaluation protocol. The
results therefore compare mechanisms under this Humanoid-Gym protocol rather
than establishing parity with or superiority over the official LCP
implementation. Appendix Table~\ref{tab:lcp-parity} enumerates this gap
explicitly: the mechanism family and coefficient anchor are aligned at the local
comparison level, but the task, robot, checkpoint/evaluation path, and hardware
evidence are not.

The primary evidence uses one robot morphology, one rough-terrain task, five
training seeds, and selected-checkpoint evaluation. The selected-vs-final
analysis shows small deltas for the soft penalty, but checkpoint dependence is
non-negligible for SC-PPO and the heuristic. The \texttt{hfield\_moderate}
replay is a diagnostic second-setting check, not a broad multi-terrain
benchmark. A task-valid MuJoCo command-shift slice
(\texttt{plane + command\_vx = 0.6}) is also available, but it remains
command-distribution evidence rather than terrain generalization. Two bounded
terrain-aware retrain probes on the main morphology also failed to produce
task-valid selected-checkpoint rows: a stair-bearing mixed rough/stairs slice
and a no-stairs \texttt{trimesh + curriculum} slice both remained
\texttt{all\_checkpoints\_collapsed}. This pattern suggests that the current
no-heights terrain-observation contract is a likely bottleneck, so the present
paper does not provide claim-grade multi-terrain evidence.

The smoothness-mechanism hyperparameters are local choices, not globally
optimized constants. The LCP coefficient is supported by a narrow
0.001/0.002/0.004 five-seed neighborhood test, and the SC-PPO threshold is
supported only as the evaluated hard-constraint setting after a local threshold
sweep. Broad hyperparameter optimality remains outside this study.

MuJoCo replay and the added torque/power proxies are sim-to-sim evidence only.
There is no real-robot validation, calibrated actuator model, system
identification, or multi-terrain/multi-robot study. The actuator-bandwidth test
is a simulation-side stress test and does not constitute hardware validation.

The multi-robot evidence is limited. Our second-morphology probe uses a single
Unitree H1 task, a single seed, and Isaac-only evaluation. These H1 results are
therefore diagnostic rather than claim-grade: they show that mechanism transfer
can differ sharply across morphologies, but they do not establish a stable
cross-morphology ranking or any hardware-transfer claim.

The mechanism explanation remains local and partly correlational. The policy
perturbation analysis gives controlled local policy-output evidence, and the
matched MuJoCo decomposition supports a control-path metric split, but neither
proves a closed-loop causal intervention from sensitivity to joint acceleration
or return. Causal identification by intervening on the policy Jacobian
independently of other training dynamics remains open.

The statistical analysis uses five seeds and should be interpreted as
descriptive uncertainty estimation rather than a large-sample hypothesis test.

External constrained-RL coverage remains incomplete. OmniSafe PPO-Lag was
tested as a limited-scope interface feasibility check and did not produce task-valid
locomotion, but this does not imply that OmniSafe, PPO-Lag, CPO, or external
constrained RL algorithms broadly fail. A local CPO-style path was also tested
through autograd/HVP, one-update, and short training-loop checks. The tensor
and one-update checks passed, and the training loop produced finite accepted
updates, but the evaluated checkpoints did not satisfy the task-validity
criteria. CPO-style training therefore remains future work rather than a primary
baseline in this study.

\section{Conclusion}

Policy-local-sensitivity regularization is most strongly supported here as a
policy-output smoothness mechanism for humanoid locomotion. In this study, a
fixed soft Jacobian penalty gives the most consistent improvements on
sensitivity, perturbation amplification, action jitter, observation-noise
robustness, and absolute actuation-effort proxies in a bounded secondary trace
analysis, while tuned reward shaping remains stronger on some downstream
metrics such as MuJoCo joint acceleration, energy per meter, and return.
Combining both mechanisms, for example by adding a lightweight action-rate term
alongside the soft Jacobian penalty or scheduling the penalty coefficient
during training, and evaluating transfer on physical hardware are natural next
steps.

\ifsetpaperbibstyle
\bibliographystyle{plainnat}
\fi
\bibliography{references}

\appendix

\section{Additional Tables}

\begin{table}[htbp]
\centering
\caption{Alignment boundary between the local LCP-style row and official LCP.}
\label{tab:lcp-parity}
\small
\resizebox{\linewidth}{!}{%
\begin{tabular}{>{\raggedright\arraybackslash}p{2.1cm}>{\raggedright\arraybackslash}p{3.3cm}>{\raggedright\arraybackslash}p{3.3cm}}
\toprule
Dimension & Aligned in the local row & Not aligned in this paper \\
\midrule
Penalty family & Soft Jacobian/Lipschitz actor penalty in PPO & No official-task or official-code equivalence claim \\
Coefficient anchor & Published fixed-coefficient anchor; local \texttt{lcp\_weight = 0.002} row & Not evidence of global hyperparameter parity \\
Actor update path & Penalty contributes directly to PPO actor updates with \texttt{create\_graph=True} & No official implementation parity claim \\
Penalty sampling & Same rollout-sample penalty family & Local row subsamples up to 64 actor observations \\
Task / terrain & Same-task comparison within this paper & Not the official LCP task stack \\
Robot / morphology & Local H1-class Humanoid-Gym row & Not the official LCP robot stack \\
Checkpoint / evaluation path & Local checkpoint sweep with Isaac/MuJoCo replay bridge & Not official checkpoints or evaluation pipeline \\
Hardware evidence & None & No real-robot evidence in this study \\
\bottomrule
\end{tabular}}
\end{table}

\begin{table}[htbp]
\centering
\caption{Local LCP coefficient sensitivity on seeds 11/17/23/29/31.}
\label{tab:lcp-weight}
\small
\resizebox{\linewidth}{!}{%
\begin{tabular}{lcccccc}
\toprule
Weight & Fall & Vel. err & Jnt acc & Jitter & Return & Sens. \\
\midrule
0.001 & 0.040 & 0.579 & 138.860 & 0.262 & 108.189 & 2.284 \\
0.002 & 0.000 & 0.490 & 117.331 & 0.212 & 118.420 & 1.890 \\
0.004 & 0.000 & 0.538 & 125.291 & 0.228 & 118.708 & 1.678 \\
\bottomrule
\end{tabular}}
\end{table}

\begin{table}[htbp]
\centering
\caption{Selected-checkpoint MuJoCo physical secondary metrics. These
trace-based proxies are supplementary sim-to-sim evidence computed on the same
five-seed selected checkpoints as Table~\ref{tab:mujoco}. Values are five-seed
means from five trace-capture episodes per selected checkpoint; per-seed rows
remain supplementary.}
\label{tab:mujoco-physical}
\small
\resizebox{\linewidth}{!}{%
\begin{tabular}{lcccc}
\toprule
Method & Torque RMS & Abs power & Abs energy & Energy/m \\
\midrule
\lcp & 24.102 & 85.659 & 877.146 & 1069.739 \\
\scppo & 26.851 & 184.312 & 1887.355 & 1129.751 \\
\heur & 25.932 & 131.616 & 1347.749 & 785.610 \\
\bottomrule
\end{tabular}}
\end{table}

\begin{table}[htbp]
\centering
\caption{Bounded single-seed Unitree H1 supplementary diagnostic. The H1 slice
uses the same metric schema as the primary morphology but is diagnostic only:
it is not claim-grade cross-morphology validation.}
\label{tab:h1-diagnostic}
\small
\resizebox{\linewidth}{!}{%
\begin{tabular}{lccccccc}
\toprule
Method & Ckpt & Fall & Vel. err & Jnt acc & Jitter & Return & Sens. \\
\midrule
Vanilla baseline & 250 & 0.000 & 0.380 & 18.881 & 0.338 & 116.250 & 14.517 \\
Revised heuristic & 200 & 0.000 & 0.382 & 15.105 & 0.275 & 115.061 & 12.997 \\
LCP-style soft penalty & 300 & 0.100 & 0.488 & 371.901 & 2.434 & 95.919 & 3.958 \\
SC-PPO 3.8 PID & 250 & 0.200 & 0.574 & 241.324 & 1.776 & 76.658 & 3.768 \\
\bottomrule
\end{tabular}}
\end{table}

\begin{table}[htbp]
\centering
\caption{Diagnostic \texttt{hfield\_moderate} MuJoCo replay using selected
checkpoints. Values are mean $\pm$ standard deviation over seeds
11/17/23/29/31. This is a supplementary second-setting terrain check rather
than a broad multi-terrain benchmark.}
\label{tab:hfield}
\small
\resizebox{\linewidth}{!}{%
\begin{tabular}{lccccc}
\toprule
Method & Fall & Vel. err & Jnt acc & Jitter & Return \\
\midrule
\lcp & $0.350\pm0.105$ & $0.832\pm0.090$ & $321.113\pm31.745$ & $0.277\pm0.050$ & $-566.367\pm123.679$ \\
\scppo & $0.500\pm0.207$ & $0.981\pm0.080$ & $386.647\pm61.783$ & $0.380\pm0.112$ & $-513.230\pm192.807$ \\
\heur & $0.400\pm0.114$ & $1.003\pm0.120$ & $322.360\pm26.072$ & $0.358\pm0.030$ & $-666.633\pm130.789$ \\
\bottomrule
\end{tabular}}
\end{table}

\begin{table}[htbp]
\centering
\caption{Task-valid MuJoCo command-shift second setting using selected
checkpoints. The replay path stays on the matched plane protocol but increases
\texttt{command\_vx} from $0.4$ to $0.6$. Values are mean $\pm$ standard
deviation over seeds 11/17/23/29/31.}
\label{tab:command-shift}
\small
\resizebox{\linewidth}{!}{%
\begin{tabular}{lcccccc}
\toprule
Method & Fall & Vel. err & Jnt acc & Jitter & Return & Sens. \\
\midrule
\heur & $0.000\pm0.000$ & $0.538\pm0.052$ & $118.274\pm17.544$ & $0.245\pm0.035$ & $-512.075\pm234.464$ & $7.331\pm0.237$ \\
\lcp & $0.000\pm0.000$ & $0.562\pm0.058$ & $119.555\pm12.496$ & $0.198\pm0.030$ & $-784.939\pm226.842$ & $1.890\pm0.064$ \\
\scppo & $0.040\pm0.080$ & $0.598\pm0.110$ & $176.182\pm87.168$ & $0.352\pm0.190$ & $-657.114\pm361.429$ & $3.630\pm0.067$ \\
\bottomrule
\end{tabular}}
\end{table}

\begin{table}[htbp]
\centering
\caption{No-retraining MuJoCo actuator-bandwidth stress. Raw jitter factors are
relative to each method's own $\tau=0$ replay. The $\tau=0.10$ setting is best
read as a task-validity boundary rather than a method-ranking regime.}
\label{tab:actuator}
\small
\resizebox{\linewidth}{!}{%
\begin{tabular}{lcccc}
\toprule
Method & $\tau=0.05$ Fall & $\tau=0.05$ Jitter fac. & $\tau=0.10$ Fall & $\tau=0.10$ Jitter fac. \\
\midrule
\lcp & $0.260$ & $1.290$ & $0.970$ & $1.888$ \\
\scppo & $0.260$ & $1.293$ & $0.960$ & $1.167$ \\
\heur & $0.330$ & $1.276$ & $0.870$ & $1.617$ \\
\bottomrule
\end{tabular}}
\end{table}

\begin{table}[htbp]
\centering
\caption{OmniSafe PPO-Lag feasibility check. Task-invalid policies are not
smooth-locomotion evidence; the near-zero jitter in the selected condition
arises because a collapsed policy barely produces action variation.}
\label{tab:omnisafe}
\small
\resizebox{\linewidth}{!}{%
\begin{tabular}{lccccc}
\toprule
Slice & Fall & Vel. err & Jnt acc & Jitter & Return \\
\midrule
Selected feasibility checkpoint & 1.000 & 1.468 & 79.309 & 0.008 & 4.386 \\
Final checkpoint 400 & 1.000 & 1.109 & 321.328 & 0.741 & 2.714 \\
\bottomrule
\end{tabular}}
\end{table}

\begin{table}[htbp]
\centering
\caption{Appendix map for bounded diagnostics and negative-result lines.}
\label{tab:appendix-map}
\small
\begin{tabular}{p{3.0cm}p{8.7cm}}
\toprule
Category & Placement and role \\
\midrule
Main-text evidence & Five-seed Isaac comparison, matched MuJoCo replay, policy-perturbation plus observation-noise mechanism diagnostics, and robustness/reliability analyses. \\
Supplementary diagnostics & MuJoCo physical proxies, \texttt{hfield\_moderate} replay, the task-valid \texttt{plane + command\_vx = 0.6} slice, actuator-bandwidth stress, the bounded Unitree H1 diagnostic, and the LCP parity table. \\
Negative-result boundaries & OmniSafe PPO-Lag, local CPO-style checks, the LCP+heuristic hybrid pilot, and the stair-bearing plus no-stairs terrain-aware retrain no-go notes remain boundary evidence rather than promoted baselines. \\
Future-work boundary & Official LCP parity, height-aware terrain-observation support for task-valid multi-terrain retraining, broader multi-robot validation, and hardware transfer stay outside the current paper. \\
\bottomrule
\end{tabular}
\end{table}

\section{Supplementary Analysis Summary}

The supplementary analyses include paired bootstrap and IQM uncertainty,
selected-vs-final checkpoint comparison, local LCP coefficient sensitivity,
policy perturbation measurements, MuJoCo physical secondary metrics,
actuator-bandwidth stress, a task-valid command-shift second setting, a bounded
Unitree H1 mechanism-transfer diagnostic, OmniSafe PPO-Lag feasibility checks,
and local CPO-style implementation checks. The physical metrics and
actuator-bandwidth stress are diagnostic simulation evidence only. The
command-shift slice is task-valid supplementary robustness evidence, but it is
not a terrain benchmark. The H1 slice is not included in the primary comparison
because it is single-seed, single-morphology supplementary evidence. The
OmniSafe and CPO-style experiments are not included in the primary comparison
because they remained limited-scope diagnostics and did not produce task-valid
locomotion under the experiments reported here. Additional supporting
negative-result notes cover the LCP+heuristic hybrid pilot, the mixed
rough/stairs retrain no-go, the no-stairs \texttt{trimesh + curriculum}
retrain no-go, and the closer-to-official full-batch LCP scout; these lines
remain claim-boundary context rather than promoted appendix tables.
The appendix map in
Table~\ref{tab:appendix-map} records the intended separation between main-text
evidence, supplementary diagnostics, negative-result boundaries, and future
work.

\section{SC-PPO Additional Comparisons}

The SC-PPO hard-constraint condition uses threshold 3.8, quantile aggregation, and
PID-Lagrangian multiplier updates. These settings define the evaluated
hard-constraint condition rather than an exhaustively optimized SC-PPO
configuration. The local threshold sweep found no shared nearby threshold that
passed all task-validity and smoothness gates, and the matched plain-dual
comparison at the same threshold also failed the task-validity criteria. These
results support the chosen PID-Lagrangian SC-PPO condition for mechanism
comparison, but they do not establish global threshold optimality or the
independent necessity of every PID term.

\end{document}
